% !TEX root = main.tex
\appendix


\clearpage
\begin{landscape}

\section{Estimation diagnostics of the Markov jump processes}
\label{sec:diagnostics_estimation}

\begin{table}[h!]
\centering
\caption{Akaike information criterion (AIC; computed from the log-likelihood at the score-optimal parameters) and per-observation gradient norm of the score objective.}
\label{tab:AIC_tab}
\scriptsize

\begin{threeparttable}
\begin{tabular}[t]{>{}lllcccccccccc}
\toprule
\multicolumn{3}{c}{ } & \multicolumn{2}{c}{Fold 1} & \multicolumn{2}{c}{Fold 2} & \multicolumn{2}{c}{Fold 3} & \multicolumn{2}{c}{Fold 4} & \multicolumn{2}{c}{Fold 5} \\
\cmidrule(l{3pt}r{3pt}){4-5} \cmidrule(l{3pt}r{3pt}){6-7} \cmidrule(l{3pt}r{3pt}){8-9} \cmidrule(l{3pt}r{3pt}){10-11} \cmidrule(l{3pt}r{3pt}){12-13}
 & Model & Exogenous info. & AIC & $\frac{1}{n}||\nabla \mathcal{D} (\bm \theta)||_2$ & AIC & $\frac{1}{n}||\nabla \mathcal{D} (\bm \theta)||_2$ & AIC & $\frac{1}{n}||\nabla \mathcal{D} (\bm \theta)||_2$ & AIC & $\frac{1}{n}||\nabla \mathcal{D} (\bm \theta)||_2$ & AIC & $\frac{1}{n}||\nabla \mathcal{D} (\bm \theta)||_2$\\
\midrule
\addlinespace[0.3em]
\multicolumn{13}{l}{\textbf{\textbf{Markov jump process}}}\\
\hspace{1em} & $\bm{A}_{\RN{1}}$: Nearest class & No & 6977.0 & 9.55e-07 & 6937.6 & 1.59e-06 & 6879.5 & 3.39e-06 & 6940.2 & 3.39e-06 & 6891.2 & 1.18e-06\\

\hspace{1em} &  & Yes (global) & 6886.6 & 3.53e-06 & 6864.3 & 3.56e-06 & 6799.7 & 5.74e-07 & 6862.7 & 3.58e-06 & 6815.0 & 5.54e-07\\

\hspace{1em} &  & Yes (state-spec.) & 6332.8 & 1.49e-06 & 6277.1 & 8.68e-07 & 6274.8 & 1.14e-06 & 6316.4 & 3.74e-07 & 6276.9 & 7.48e-07\\

\addlinespace[0.5em]

\hspace{1em} & $\bm{A}_{\RN{2}}$: Relaxed Erlang & No & 6396.8 & 9.47e-07 & 6327.1 & 2.53e-06 & 6298.3 & 7.17e-07 & 6349.0 & 1.38e-06 & 6321.6 & 1.35e-06\\

\hspace{1em} &  & Yes (global) & 6322.3 & 1.79e-06 & 6270.5 & 5.78e-07 & 6233.5 & 6.22e-07 & 6275.2 & 3.69e-06 & 6260.8 & 2.25e-06\\

\hspace{1em} &  & Yes (state-spec.) & 6235.5 & 5.44e-07 & 6199.6 & 1.17e-06 & 6159.8 & 1.37e-06 & 6210.0 & 5.87e-07 & 6182.6 & 2.07e-07\\

\addlinespace[0.5em]

\hspace{1em} & $\bm{A}_{\RN{3}}$: Up to 2 classes ahead & No & 6159.5 & 2.79e-06 & 6108.7 & 1.51e-07 & 6085.4 & 1.74e-06 & 6135.0 & 2.08e-06 & 6093.6 & 4.66e-07\\

\hspace{1em} &  & Yes (global) & 6074.7 & 3.01e-06 & 6038.8 & 1.40e-06 & 6004.0 & 1.75e-06 & 6063.0 & 1.32e-06 & 6023.4 & 1.83e-06\\

\hspace{1em} &  & Yes (state-spec.) & 5969.2 & 2.52e-07 & 5931.1 & 3.30e-07 & 5898.0 & 8.74e-07 & 5952.3 & 5.96e-07 & 5912.5 & 5.99e-07\\

\addlinespace[0.5em]

\hspace{1em} & $\bm{A}_{\RN{4}}$: Up to 3 classes ahead & No & 6125.1 & 1.05e-07 & 6053.4 & 2.95e-06 & 6043.8 & 1.54e-06 & 6093.9 & 2.29e-06 & 6057.1 & 1.87e-06\\

\hspace{1em} &  & Yes (global) & 6039.1 & 1.99e-06 & 5981.9 & 7.29e-07 & 5962.9 & 3.12e-06 & 6017.3 & 3.02e-06 & 5986.3 & 4.22e-08\\

\hspace{1em} &  & Yes (state-spec.) & 5971.5 & 2.79e-07 & 5933.4 & 5.83e-07 & 5897.5 & 3.76e-07 & 5954.1 & 4.58e-07 & 5915.2 & 1.77e-06\\

\addlinespace[0.5em]

\hspace{1em} & $\bm{A}_{\RN{5}}$: Free upper-triangular & No & 6124.8 & 8.66e-07 & 6054.7 & 3.87e-08 & 6043.2 & 3.44e-08 & 6093.9 & 6.81e-08 & 6059.1 & 2.76e-09\\

\hspace{1em} &  & Yes (global) & 6039.5 & 3.46e-07 & 5983.7 & 4.77e-07 & 5963.1 & 1.44e-06 & 6017.4 & 9.82e-07 & 5988.3 & 1.23e-08\\

\hspace{1em} &  & Yes (state-spec.) & 5973.5 & 1.23e-06 & 5935.4 & 7.12e-07 & 5899.3 & 6.27e-07 & 5956.1 & 3.74e-07 & 5917.2 & 1.42e-06\\
\cmidrule{1-13}
\addlinespace[0.3em]
\multicolumn{13}{l}{\textbf{\textbf{\makecell[l]{\textbf{Model with modified}\\ \textbf{sojourn time estimates}}}}}\\
\hspace{1em} & $\bm{A}_{\RN{1}}$: Nearest class & No & 6620.3 & 1.77e-06 & 6523.1 & 2.72e-06 & 6549.4 & 3.07e-06 & 6588.9 & 7.48e-07 & 6539.3 & 1.87e-06\\

\hspace{1em} &  & Yes (global) & 6500.8 & 2.07e-06 & 6423.4 & 1.24e-06 & 6438.9 & 1.02e-06 & 6483.5 & 1.53e-06 & 6436.7 & 2.32e-06\\

\addlinespace[0.5em]

\hspace{1em} & $\bm{A}_{\RN{2}}$: Relaxed Erlang & No & 6378.1 & 2.17e-06 & 6301.1 & 3.17e-06 & 6285.2 & 9.20e-07 & 6338.7 & 1.52e-06 & 6308.4 & 1.70e-06\\

\hspace{1em} &  & Yes (global) & 6301.7 & 1.38e-06 & 6241.3 & 4.86e-06 & 6217.2 & 2.03e-06 & 6265.5 & 2.36e-06 & 6247.9 & 1.40e-06\\

\addlinespace[0.5em]

\hspace{1em} & $\bm{A}_{\RN{3}}$: Up to 2 classes ahead & No & 6131.4 & 2.81e-06 & 6063.9 & 3.36e-06 & 6058.7 & 2.02e-06 & 6102.6 & 2.64e-06 & 6066.8 & 1.18e-06\\

\hspace{1em} &  & Yes (global) & 6044.8 & 1.34e-06 & 5992.0 & 2.07e-06 & 5976.8 & 2.21e-06 & 6032.1 & 1.21e-06 & 5994.5 & 1.57e-06\\

\addlinespace[0.5em]

\hspace{1em} & $\bm{A}_{\RN{4}}$: Up to 3 classes ahead & No & 6122.2 & 5.50e-07 & 6050.2 & 8.50e-07 & 6042.8 & 4.70e-07 & 6089.9 & 2.14e-06 & 6055.5 & 8.03e-07\\

\hspace{1em} &  & Yes (global) & 6037.9 & 9.81e-07 & 5980.5 & 1.26e-06 & 5963.8 & 1.50e-06 & 6019.6 & 2.69e-06 & 5986.7 & 1.79e-06\\

\addlinespace[0.5em]

\hspace{1em} & $\bm{A}_{\RN{5}}$: Free upper-triangular & No & 6122.8 & 1.54e-06 & 6051.9 & 1.70e-06 & 6043.2 & 1.20e-06 & 6090.8 & 1.40e-06 & 6057.5 & 4.96e-08\\

\hspace{1em} &  & Yes (global) & 6039.2 & 6.09e-07 & 5982.5 & 1.77e-07 & 5964.9 & 2.36e-06 & 6020.8 & 6.55e-07 & 5988.7 & 4.90e-08\\
\cmidrule{1-13}
\addlinespace[0.3em]
\multicolumn{13}{l}{\textbf{\textbf{Mixture Markov jump process ($G=4$)}}}\\
\hspace{1em} & $\bm{A}_{\RN{1}}$: Nearest class & No & 5804.7 & 5.45e-07 & 6967.6 & 2.18e-06 & 5789.8 & 8.88e-07 & 6970.2 & 7.67e-07 & 5770.5 & 8.39e-07\\

\hspace{1em} &  & Yes (global) & 5664.3 & 1.28e-06 & 5641.5 & 8.96e-04 & 5673.5 & 9.71e-07 & 5702.9 & 1.40e-07 & 5677.4 & 1.30e-06\\

\addlinespace[0.5em]

\hspace{1em} & $\bm{A}_{\RN{2}}$: Relaxed Erlang & No & 6426.8 & 1.32e-06 & 6357.1 & 1.62e-07 & 6328.3 & 1.48e-06 & 6379.0 & 1.44e-06 & 6351.6 & 3.88e-07\\

\hspace{1em} &  & Yes (global) & 5678.9 & 1.11e-07 & 5639.1 & 8.72e-08 & 5662.6 & 7.79e-07 & 5643.5 & 9.30e-07 & 5673.2 & 1.19e-04\\

\addlinespace[0.5em]

\hspace{1em} & $\bm{A}_{\RN{3}}$: Up to 2 classes ahead & No & 5695.0 & 6.53e-07 & 5654.8 & 1.01e-06 & 5668.8 & 4.59e-07 & 5707.5 & 4.41e-07 & 5685.5 & 6.38e-07\\

\hspace{1em} &  & Yes (global) & 5606.3 & 3.74e-07 & 5583.3 & 5.15e-07 & 5578.1 & 1.03e-06 & 5607.3 & 7.46e-07 & 5607.3 & 3.59e-07\\

\addlinespace[0.5em]

\hspace{1em} & $\bm{A}_{\RN{4}}$: Up to 3 classes ahead & No & 5758.8 & 6.09e-05 & 5667.4 & 6.03e-07 & 5670.3 & 7.92e-07 & 5694.8 & 9.54e-07 & 5686.9 & 5.11e-07\\

\hspace{1em} &  & Yes (global) & 5621.9 & 8.96e-07 & 5595.7 & 8.78e-07 & 5593.9 & 5.27e-07 & 5627.6 & 7.94e-07 & 5619.9 & 7.16e-07\\

\addlinespace[0.5em]

\hspace{1em} & $\bm{A}_{\RN{5}}$: Free upper-triangular & No & 5703.9 & 1.50e-06 & 5677.6 & 1.15e-06 & 5675.9 & 1.67e-06 & 5704.0 & 5.47e-07 & 5694.4 & 8.27e-07\\

\hspace{1em} &  & Yes (global) & 5630.0 & 9.15e-07 & 5596.8 & 8.01e-07 & 5596.4 & 3.61e-07 & 5635.7 & 2.48e-06 & 5624.8 & 1.10e-06\\
\bottomrule
\end{tabular}
\begin{tablenotes}
\item[1] Exogenous information indicates whether the rail characteristics described in \Cref{sec:data} are included: Yes (global) uses \cref{eq:trans_rates_covariates}; Yes (state-spec.) uses \cref{eq:trans_rates_covariates2}.
\end{tablenotes}
\end{threeparttable}

\end{table}

\end{landscape}
\clearpage

%\section{Estimation diagnostics of the Markov jump processes} \label{sec:diagnostics_estimation}
%
%\begin{table}[h!]  
%\centering
%\caption{Akaike information criterion (AIC; computed from the log-likelihood at the score-optimal parameters) and per-observation gradient norm of the score objective, reported by fold in the cross-validation.} \label{tab:AIC_tab}
%\centering
%\scriptsize
%
\begin{threeparttable}
\begin{tabular}[t]{>{}lllcccccccccc}
\toprule
\multicolumn{3}{c}{ } & \multicolumn{2}{c}{Fold 1} & \multicolumn{2}{c}{Fold 2} & \multicolumn{2}{c}{Fold 3} & \multicolumn{2}{c}{Fold 4} & \multicolumn{2}{c}{Fold 5} \\
\cmidrule(l{3pt}r{3pt}){4-5} \cmidrule(l{3pt}r{3pt}){6-7} \cmidrule(l{3pt}r{3pt}){8-9} \cmidrule(l{3pt}r{3pt}){10-11} \cmidrule(l{3pt}r{3pt}){12-13}
 & Model & Exogenous info. & AIC & $\frac{1}{n}||\nabla \mathcal{D} (\bm \theta)||_2$ & AIC & $\frac{1}{n}||\nabla \mathcal{D} (\bm \theta)||_2$ & AIC & $\frac{1}{n}||\nabla \mathcal{D} (\bm \theta)||_2$ & AIC & $\frac{1}{n}||\nabla \mathcal{D} (\bm \theta)||_2$ & AIC & $\frac{1}{n}||\nabla \mathcal{D} (\bm \theta)||_2$\\
\midrule
\addlinespace[0.3em]
\multicolumn{13}{l}{\textbf{\textbf{Markov jump process}}}\\
\hspace{1em} & $\bm{A}_{\RN{1}}$: Nearest class & No & 6977.0 & 9.55e-07 & 6937.6 & 1.59e-06 & 6879.5 & 3.39e-06 & 6940.2 & 3.39e-06 & 6891.2 & 1.18e-06\\

\hspace{1em} &  & Yes (global) & 6886.6 & 3.53e-06 & 6864.3 & 3.56e-06 & 6799.7 & 5.74e-07 & 6862.7 & 3.58e-06 & 6815.0 & 5.54e-07\\

\hspace{1em} &  & Yes (state-spec.) & 6332.8 & 1.49e-06 & 6277.1 & 8.68e-07 & 6274.8 & 1.14e-06 & 6316.4 & 3.74e-07 & 6276.9 & 7.48e-07\\

\addlinespace[0.5em]

\hspace{1em} & $\bm{A}_{\RN{2}}$: Relaxed Erlang & No & 6396.8 & 9.47e-07 & 6327.1 & 2.53e-06 & 6298.3 & 7.17e-07 & 6349.0 & 1.38e-06 & 6321.6 & 1.35e-06\\

\hspace{1em} &  & Yes (global) & 6322.3 & 1.79e-06 & 6270.5 & 5.78e-07 & 6233.5 & 6.22e-07 & 6275.2 & 3.69e-06 & 6260.8 & 2.25e-06\\

\hspace{1em} &  & Yes (state-spec.) & 6235.5 & 5.44e-07 & 6199.6 & 1.17e-06 & 6159.8 & 1.37e-06 & 6210.0 & 5.87e-07 & 6182.6 & 2.07e-07\\

\addlinespace[0.5em]

\hspace{1em} & $\bm{A}_{\RN{3}}$: Up to 2 classes ahead & No & 6159.5 & 2.79e-06 & 6108.7 & 1.51e-07 & 6085.4 & 1.74e-06 & 6135.0 & 2.08e-06 & 6093.6 & 4.66e-07\\

\hspace{1em} &  & Yes (global) & 6074.7 & 3.01e-06 & 6038.8 & 1.40e-06 & 6004.0 & 1.75e-06 & 6063.0 & 1.32e-06 & 6023.4 & 1.83e-06\\

\hspace{1em} &  & Yes (state-spec.) & 5969.2 & 2.52e-07 & 5931.1 & 3.30e-07 & 5898.0 & 8.74e-07 & 5952.3 & 5.96e-07 & 5912.5 & 5.99e-07\\

\addlinespace[0.5em]

\hspace{1em} & $\bm{A}_{\RN{4}}$: Up to 3 classes ahead & No & 6125.1 & 1.05e-07 & 6053.4 & 2.95e-06 & 6043.8 & 1.54e-06 & 6093.9 & 2.29e-06 & 6057.1 & 1.87e-06\\

\hspace{1em} &  & Yes (global) & 6039.1 & 1.99e-06 & 5981.9 & 7.29e-07 & 5962.9 & 3.12e-06 & 6017.3 & 3.02e-06 & 5986.3 & 4.22e-08\\

\hspace{1em} &  & Yes (state-spec.) & 5971.5 & 2.79e-07 & 5933.4 & 5.83e-07 & 5897.5 & 3.76e-07 & 5954.1 & 4.58e-07 & 5915.2 & 1.77e-06\\

\addlinespace[0.5em]

\hspace{1em} & $\bm{A}_{\RN{5}}$: Free upper-triangular & No & 6124.8 & 8.66e-07 & 6054.7 & 3.87e-08 & 6043.2 & 3.44e-08 & 6093.9 & 6.81e-08 & 6059.1 & 2.76e-09\\

\hspace{1em} &  & Yes (global) & 6039.5 & 3.46e-07 & 5983.7 & 4.77e-07 & 5963.1 & 1.44e-06 & 6017.4 & 9.82e-07 & 5988.3 & 1.23e-08\\

\hspace{1em} &  & Yes (state-spec.) & 5973.5 & 1.23e-06 & 5935.4 & 7.12e-07 & 5899.3 & 6.27e-07 & 5956.1 & 3.74e-07 & 5917.2 & 1.42e-06\\
\cmidrule{1-13}
\addlinespace[0.3em]
\multicolumn{13}{l}{\textbf{\textbf{\makecell[l]{\textbf{Model with modified}\\ \textbf{sojourn time estimates}}}}}\\
\hspace{1em} & $\bm{A}_{\RN{1}}$: Nearest class & No & 6620.3 & 1.77e-06 & 6523.1 & 2.72e-06 & 6549.4 & 3.07e-06 & 6588.9 & 7.48e-07 & 6539.3 & 1.87e-06\\

\hspace{1em} &  & Yes (global) & 6500.8 & 2.07e-06 & 6423.4 & 1.24e-06 & 6438.9 & 1.02e-06 & 6483.5 & 1.53e-06 & 6436.7 & 2.32e-06\\

\addlinespace[0.5em]

\hspace{1em} & $\bm{A}_{\RN{2}}$: Relaxed Erlang & No & 6378.1 & 2.17e-06 & 6301.1 & 3.17e-06 & 6285.2 & 9.20e-07 & 6338.7 & 1.52e-06 & 6308.4 & 1.70e-06\\

\hspace{1em} &  & Yes (global) & 6301.7 & 1.38e-06 & 6241.3 & 4.86e-06 & 6217.2 & 2.03e-06 & 6265.5 & 2.36e-06 & 6247.9 & 1.40e-06\\

\addlinespace[0.5em]

\hspace{1em} & $\bm{A}_{\RN{3}}$: Up to 2 classes ahead & No & 6131.4 & 2.81e-06 & 6063.9 & 3.36e-06 & 6058.7 & 2.02e-06 & 6102.6 & 2.64e-06 & 6066.8 & 1.18e-06\\

\hspace{1em} &  & Yes (global) & 6044.8 & 1.34e-06 & 5992.0 & 2.07e-06 & 5976.8 & 2.21e-06 & 6032.1 & 1.21e-06 & 5994.5 & 1.57e-06\\

\addlinespace[0.5em]

\hspace{1em} & $\bm{A}_{\RN{4}}$: Up to 3 classes ahead & No & 6122.2 & 5.50e-07 & 6050.2 & 8.50e-07 & 6042.8 & 4.70e-07 & 6089.9 & 2.14e-06 & 6055.5 & 8.03e-07\\

\hspace{1em} &  & Yes (global) & 6037.9 & 9.81e-07 & 5980.5 & 1.26e-06 & 5963.8 & 1.50e-06 & 6019.6 & 2.69e-06 & 5986.7 & 1.79e-06\\

\addlinespace[0.5em]

\hspace{1em} & $\bm{A}_{\RN{5}}$: Free upper-triangular & No & 6122.8 & 1.54e-06 & 6051.9 & 1.70e-06 & 6043.2 & 1.20e-06 & 6090.8 & 1.40e-06 & 6057.5 & 4.96e-08\\

\hspace{1em} &  & Yes (global) & 6039.2 & 6.09e-07 & 5982.5 & 1.77e-07 & 5964.9 & 2.36e-06 & 6020.8 & 6.55e-07 & 5988.7 & 4.90e-08\\
\cmidrule{1-13}
\addlinespace[0.3em]
\multicolumn{13}{l}{\textbf{\textbf{Mixture Markov jump process ($G=4$)}}}\\
\hspace{1em} & $\bm{A}_{\RN{1}}$: Nearest class & No & 5804.7 & 5.45e-07 & 6967.6 & 2.18e-06 & 5789.8 & 8.88e-07 & 6970.2 & 7.67e-07 & 5770.5 & 8.39e-07\\

\hspace{1em} &  & Yes (global) & 5664.3 & 1.28e-06 & 5641.5 & 8.96e-04 & 5673.5 & 9.71e-07 & 5702.9 & 1.40e-07 & 5677.4 & 1.30e-06\\

\addlinespace[0.5em]

\hspace{1em} & $\bm{A}_{\RN{2}}$: Relaxed Erlang & No & 6426.8 & 1.32e-06 & 6357.1 & 1.62e-07 & 6328.3 & 1.48e-06 & 6379.0 & 1.44e-06 & 6351.6 & 3.88e-07\\

\hspace{1em} &  & Yes (global) & 5678.9 & 1.11e-07 & 5639.1 & 8.72e-08 & 5662.6 & 7.79e-07 & 5643.5 & 9.30e-07 & 5673.2 & 1.19e-04\\

\addlinespace[0.5em]

\hspace{1em} & $\bm{A}_{\RN{3}}$: Up to 2 classes ahead & No & 5695.0 & 6.53e-07 & 5654.8 & 1.01e-06 & 5668.8 & 4.59e-07 & 5707.5 & 4.41e-07 & 5685.5 & 6.38e-07\\

\hspace{1em} &  & Yes (global) & 5606.3 & 3.74e-07 & 5583.3 & 5.15e-07 & 5578.1 & 1.03e-06 & 5607.3 & 7.46e-07 & 5607.3 & 3.59e-07\\

\addlinespace[0.5em]

\hspace{1em} & $\bm{A}_{\RN{4}}$: Up to 3 classes ahead & No & 5758.8 & 6.09e-05 & 5667.4 & 6.03e-07 & 5670.3 & 7.92e-07 & 5694.8 & 9.54e-07 & 5686.9 & 5.11e-07\\

\hspace{1em} &  & Yes (global) & 5621.9 & 8.96e-07 & 5595.7 & 8.78e-07 & 5593.9 & 5.27e-07 & 5627.6 & 7.94e-07 & 5619.9 & 7.16e-07\\

\addlinespace[0.5em]

\hspace{1em} & $\bm{A}_{\RN{5}}$: Free upper-triangular & No & 5703.9 & 1.50e-06 & 5677.6 & 1.15e-06 & 5675.9 & 1.67e-06 & 5704.0 & 5.47e-07 & 5694.4 & 8.27e-07\\

\hspace{1em} &  & Yes (global) & 5630.0 & 9.15e-07 & 5596.8 & 8.01e-07 & 5596.4 & 3.61e-07 & 5635.7 & 2.48e-06 & 5624.8 & 1.10e-06\\
\bottomrule
\end{tabular}
\begin{tablenotes}
\item[1] Exogenous information indicates whether the rail characteristics described in \Cref{sec:data} are included: Yes (global) uses \cref{eq:trans_rates_covariates}; Yes (state-spec.) uses \cref{eq:trans_rates_covariates2}.
\end{tablenotes}
\end{threeparttable}

%%    \resizebox{\ifdim\width>\linewidth\linewidth\else\width\fi}{!}{
\begin{threeparttable}
\begin{tabular}[t]{>{}lllcccccccccc}
\toprule
\multicolumn{3}{c}{ } & \multicolumn{2}{c}{Fold 1} & \multicolumn{2}{c}{Fold 2} & \multicolumn{2}{c}{Fold 3} & \multicolumn{2}{c}{Fold 4} & \multicolumn{2}{c}{Fold 5} \\
\cmidrule(l{3pt}r{3pt}){4-5} \cmidrule(l{3pt}r{3pt}){6-7} \cmidrule(l{3pt}r{3pt}){8-9} \cmidrule(l{3pt}r{3pt}){10-11} \cmidrule(l{3pt}r{3pt}){12-13}
 & Model & Exogenous info. & AIC & $\frac{1}{n}||\nabla \mathcal{D} (\bm \theta)||_2$ & AIC & $\frac{1}{n}||\nabla \mathcal{D} (\bm \theta)||_2$ & AIC & $\frac{1}{n}||\nabla \mathcal{D} (\bm \theta)||_2$ & AIC & $\frac{1}{n}||\nabla \mathcal{D} (\bm \theta)||_2$ & AIC & $\frac{1}{n}||\nabla \mathcal{D} (\bm \theta)||_2$\\
\midrule
\addlinespace[0.3em]
\multicolumn{13}{l}{\textbf{\textbf{Markov jump process}}}\\
\hspace{1em} & $\bm{A}_{\RN{1}}$: Nearest class & No & 6977.0 & 9.55e-07 & 6937.6 & 1.59e-06 & 6879.5 & 3.39e-06 & 6940.2 & 3.39e-06 & 6891.2 & 1.18e-06\\

\hspace{1em} &  & Yes (global) & 6886.6 & 3.53e-06 & 6864.3 & 3.56e-06 & 6799.7 & 5.74e-07 & 6862.7 & 3.58e-06 & 6815.0 & 5.54e-07\\

\hspace{1em} &  & Yes (state-spec.) & 6332.8 & 1.49e-06 & 6277.1 & 8.68e-07 & 6274.8 & 1.14e-06 & 6316.4 & 3.74e-07 & 6276.9 & 7.48e-07\\

\addlinespace[0.5em]

\hspace{1em} & $\bm{A}_{\RN{2}}$: Relaxed Erlang & No & 6396.8 & 9.47e-07 & 6327.1 & 2.53e-06 & 6298.3 & 7.17e-07 & 6349.0 & 1.38e-06 & 6321.6 & 1.35e-06\\

\hspace{1em} &  & Yes (global) & 6322.3 & 1.79e-06 & 6270.5 & 5.78e-07 & 6233.5 & 6.22e-07 & 6275.2 & 3.69e-06 & 6260.8 & 2.25e-06\\

\hspace{1em} &  & Yes (state-spec.) & 6235.5 & 5.44e-07 & 6199.6 & 1.17e-06 & 6159.8 & 1.37e-06 & 6210.0 & 5.87e-07 & 6182.6 & 2.07e-07\\

\addlinespace[0.5em]

\hspace{1em} & $\bm{A}_{\RN{3}}$: Up to 2 classes ahead & No & 6159.5 & 2.79e-06 & 6108.7 & 1.51e-07 & 6085.4 & 1.74e-06 & 6135.0 & 2.08e-06 & 6093.6 & 4.66e-07\\

\hspace{1em} &  & Yes (global) & 6074.7 & 3.01e-06 & 6038.8 & 1.40e-06 & 6004.0 & 1.75e-06 & 6063.0 & 1.32e-06 & 6023.4 & 1.83e-06\\

\hspace{1em} &  & Yes (state-spec.) & 5969.2 & 2.52e-07 & 5931.1 & 3.30e-07 & 5898.0 & 8.74e-07 & 5952.3 & 5.96e-07 & 5912.5 & 5.99e-07\\

\addlinespace[0.5em]

\hspace{1em} & $\bm{A}_{\RN{4}}$: Up to 3 classes ahead & No & 6125.1 & 1.05e-07 & 6053.4 & 2.95e-06 & 6043.8 & 1.54e-06 & 6093.9 & 2.29e-06 & 6057.1 & 1.87e-06\\

\hspace{1em} &  & Yes (global) & 6039.1 & 1.99e-06 & 5981.9 & 7.29e-07 & 5962.9 & 3.12e-06 & 6017.3 & 3.02e-06 & 5986.3 & 4.22e-08\\

\hspace{1em} &  & Yes (state-spec.) & 5971.5 & 2.79e-07 & 5933.4 & 5.83e-07 & 5897.5 & 3.76e-07 & 5954.1 & 4.58e-07 & 5915.2 & 1.77e-06\\

\addlinespace[0.5em]

\hspace{1em} & $\bm{A}_{\RN{5}}$: Free upper-triangular & No & 6124.8 & 8.66e-07 & 6054.7 & 3.87e-08 & 6043.2 & 3.44e-08 & 6093.9 & 6.81e-08 & 6059.1 & 2.76e-09\\

\hspace{1em} &  & Yes (global) & 6039.5 & 3.46e-07 & 5983.7 & 4.77e-07 & 5963.1 & 1.44e-06 & 6017.4 & 9.82e-07 & 5988.3 & 1.23e-08\\

\hspace{1em} &  & Yes (state-spec.) & 5973.5 & 1.23e-06 & 5935.4 & 7.12e-07 & 5899.3 & 6.27e-07 & 5956.1 & 3.74e-07 & 5917.2 & 1.42e-06\\
\cmidrule{1-13}
\addlinespace[0.3em]
\multicolumn{13}{l}{\textbf{\textbf{\makecell[l]{\textbf{Model with modified}\\ \textbf{sojourn time estimates}}}}}\\
\hspace{1em} & $\bm{A}_{\RN{1}}$: Nearest class & No & 6620.3 & 1.77e-06 & 6523.1 & 2.72e-06 & 6549.4 & 3.07e-06 & 6588.9 & 7.48e-07 & 6539.3 & 1.87e-06\\

\hspace{1em} &  & Yes (global) & 6500.8 & 2.07e-06 & 6423.4 & 1.24e-06 & 6438.9 & 1.02e-06 & 6483.5 & 1.53e-06 & 6436.7 & 2.32e-06\\

\addlinespace[0.5em]

\hspace{1em} & $\bm{A}_{\RN{2}}$: Relaxed Erlang & No & 6378.1 & 2.17e-06 & 6301.1 & 3.17e-06 & 6285.2 & 9.20e-07 & 6338.7 & 1.52e-06 & 6308.4 & 1.70e-06\\

\hspace{1em} &  & Yes (global) & 6301.7 & 1.38e-06 & 6241.3 & 4.86e-06 & 6217.2 & 2.03e-06 & 6265.5 & 2.36e-06 & 6247.9 & 1.40e-06\\

\addlinespace[0.5em]

\hspace{1em} & $\bm{A}_{\RN{3}}$: Up to 2 classes ahead & No & 6131.4 & 2.81e-06 & 6063.9 & 3.36e-06 & 6058.7 & 2.02e-06 & 6102.6 & 2.64e-06 & 6066.8 & 1.18e-06\\

\hspace{1em} &  & Yes (global) & 6044.8 & 1.34e-06 & 5992.0 & 2.07e-06 & 5976.8 & 2.21e-06 & 6032.1 & 1.21e-06 & 5994.5 & 1.57e-06\\

\addlinespace[0.5em]

\hspace{1em} & $\bm{A}_{\RN{4}}$: Up to 3 classes ahead & No & 6122.2 & 5.50e-07 & 6050.2 & 8.50e-07 & 6042.8 & 4.70e-07 & 6089.9 & 2.14e-06 & 6055.5 & 8.03e-07\\

\hspace{1em} &  & Yes (global) & 6037.9 & 9.81e-07 & 5980.5 & 1.26e-06 & 5963.8 & 1.50e-06 & 6019.6 & 2.69e-06 & 5986.7 & 1.79e-06\\

\addlinespace[0.5em]

\hspace{1em} & $\bm{A}_{\RN{5}}$: Free upper-triangular & No & 6122.8 & 1.54e-06 & 6051.9 & 1.70e-06 & 6043.2 & 1.20e-06 & 6090.8 & 1.40e-06 & 6057.5 & 4.96e-08\\

\hspace{1em} &  & Yes (global) & 6039.2 & 6.09e-07 & 5982.5 & 1.77e-07 & 5964.9 & 2.36e-06 & 6020.8 & 6.55e-07 & 5988.7 & 4.90e-08\\
\cmidrule{1-13}
\addlinespace[0.3em]
\multicolumn{13}{l}{\textbf{\textbf{Mixture Markov jump process ($G=4$)}}}\\
\hspace{1em} & $\bm{A}_{\RN{1}}$: Nearest class & No & 5804.7 & 5.45e-07 & 6967.6 & 2.18e-06 & 5789.8 & 8.88e-07 & 6970.2 & 7.67e-07 & 5770.5 & 8.39e-07\\

\hspace{1em} &  & Yes (global) & 5664.3 & 1.28e-06 & 5641.5 & 8.96e-04 & 5673.5 & 9.71e-07 & 5702.9 & 1.40e-07 & 5677.4 & 1.30e-06\\

\addlinespace[0.5em]

\hspace{1em} & $\bm{A}_{\RN{2}}$: Relaxed Erlang & No & 6426.8 & 1.32e-06 & 6357.1 & 1.62e-07 & 6328.3 & 1.48e-06 & 6379.0 & 1.44e-06 & 6351.6 & 3.88e-07\\

\hspace{1em} &  & Yes (global) & 5678.9 & 1.11e-07 & 5639.1 & 8.72e-08 & 5662.6 & 7.79e-07 & 5643.5 & 9.30e-07 & 5673.2 & 1.19e-04\\

\addlinespace[0.5em]

\hspace{1em} & $\bm{A}_{\RN{3}}$: Up to 2 classes ahead & No & 5695.0 & 6.53e-07 & 5654.8 & 1.01e-06 & 5668.8 & 4.59e-07 & 5707.5 & 4.41e-07 & 5685.5 & 6.38e-07\\

\hspace{1em} &  & Yes (global) & 5606.3 & 3.74e-07 & 5583.3 & 5.15e-07 & 5578.1 & 1.03e-06 & 5607.3 & 7.46e-07 & 5607.3 & 3.59e-07\\

\addlinespace[0.5em]

\hspace{1em} & $\bm{A}_{\RN{4}}$: Up to 3 classes ahead & No & 5758.8 & 6.09e-05 & 5667.4 & 6.03e-07 & 5670.3 & 7.92e-07 & 5694.8 & 9.54e-07 & 5686.9 & 5.11e-07\\

\hspace{1em} &  & Yes (global) & 5621.9 & 8.96e-07 & 5595.7 & 8.78e-07 & 5593.9 & 5.27e-07 & 5627.6 & 7.94e-07 & 5619.9 & 7.16e-07\\

\addlinespace[0.5em]

\hspace{1em} & $\bm{A}_{\RN{5}}$: Free upper-triangular & No & 5703.9 & 1.50e-06 & 5677.6 & 1.15e-06 & 5675.9 & 1.67e-06 & 5704.0 & 5.47e-07 & 5694.4 & 8.27e-07\\

\hspace{1em} &  & Yes (global) & 5630.0 & 9.15e-07 & 5596.8 & 8.01e-07 & 5596.4 & 3.61e-07 & 5635.7 & 2.48e-06 & 5624.8 & 1.10e-06\\
\bottomrule
\end{tabular}
\begin{tablenotes}
\item[1] Exogenous information indicates whether the rail characteristics described in \Cref{sec:data} are included: Yes (global) uses \cref{eq:trans_rates_covariates}; Yes (state-spec.) uses \cref{eq:trans_rates_covariates2}.
\end{tablenotes}
\end{threeparttable}
}
%\end{table} 

\newpage

\section{Transition residuals and sharpness through entropy} \label{sec:transition_L1}  \label{sec:sharpness}
Transition residual matrices  $\bm R$ are computed using held-out predictions from the cross-validation procedure in \cref{sec:prediction_error}.
We report the L1-norm of the residual matrix obtained by pooling observed transition counts and pooled model-implied expected counts across folds, i.e. $|| \bm R || = \sum_{i,j} |  R_{i,j}|$ for $i,j \in \mathcal J$. %Lower values indicate closer agreement between observed and model-implied transition frequencies.
Further, we report mean predictive entropy as a measure of sharpness,
$H(\widehat{\boldsymbol\pi}^{(s)}) = - \frac{1}{\log m}\sum_{j=1}^{m}\widehat{\pi}^{(s)}_j \log\!\big(\widehat{\pi}^{(s)}_j\big)$,
computed for each held-out forecast and averaged. Lower entropy indicates sharper predictive distributions.

\begin{table}[h!]  
\centering
\caption{L1-norm of transition residuals $
|| \bm R || $ and normalized predictive entropy $H$ for different model settings.} \label{tab:transition_structure_L1}
\centering
\scriptsize
   \input{transition_structure_L1}
\end{table} 



\newpage

\section{Synthetic data} \label{sec:synthetic_data}
As the raw defect data provided by the infrastructure manager are confidential, they cannot be shared. To support reproducibility of the proposed methods of this study, we provide a synthetic dataset in the same format (with censoring) as the original transition data, i.e., tuples $\{(y_1^{(s)},y_2^{(s)},\tau^{(s)},\bm z^{(s)})\}_{s \in \mathcal S}$, which can be used to run all proposed estimation and prediction procedures.

The synthetic data is generated as follows.
\begin{itemize}
\item An MJP with a chosen model specification is estimated on the confidential defect data using the methods, so that model parameters $\widehat{\bm \theta}$ are obtained.
\item $n$ initial states $y_1'$ are sampled independently from the empirical distribution of the observed initial states in the confidential dataset, i.e.
$
y_1' \sim \text{Categorical}(\widehat{\bm p})$ where $\widehat p_i = \frac{1}{| \mathcal S|} \sum_{s\in\mathcal S} \mathds{1}_{y_1^{(s)}=\;i}
$ for $i \in \mathcal J$. 
%$\left(\widehat p_i = \frac{1}{|\mathcal S|}\sum_{s\in\mathcal S}\mathds{1}_{y_1^{(s)}=i}\right)$
\item $n$ observation intervals $\tau'_1,\dots,\tau'_n$ are sampled independently by bootstrap resampling from the observed intervals, i.e.,
$
\tau'_r \overset{\text{i.i.d.}}{\sim} \{\tau^{(s)}: s\in\mathcal S\} $ with replacement for $ r=1,\dots,n$. The same is done for the $p$ covariates, so that also $n$ observations are sampled marginally for each variable, i.e., $(z'_{1,1},\dots,z'_{1,n} ), \dots, (z'_{p,1},\dots,z'_{p,n})$. 
\item The transient distribution for each sampled defect under the model and the estimated parameters is then used as the categorical distribution for which an ending state is sampled, i.e., 
$
y_{2,r}' \sim \text{Categorical}(\widehat{\bm\pi}( \tau'_r \mid y'_{1,r}, \bm z'_r; \widehat{\bm\theta}) )
$ for $r = 1, \dots, n$, with $\bm z'_r = (z'_{1,r},\dots,z'_{p,r})^\top$.
\end{itemize}
We generate a synthetic dataset using the $\bm{A}_{\RN{3}}$ parameterization (\cref{sec:generator_params}), warping (\cref{sec:warping}), and global regression (\cref{eq:trans_rates_covariates} in \cref{sec:regression}) with $n = 10000$.

This synthetic dataset serves two purposes: 1) the proposed methods can be tested on data with the same structure as the confidential dataset; 2) it enables numerical verification that the estimation procedure can recover the data-generating parameters and that estimation error decreases as $n$ increases.




\newpage

\section{Details on the stochastic optimization program} \label{sec:stochastic_optimization_details}
Here, we provide additional details on the stochastic optimization program described in \cref{sec:proposition_for_use}.
\paragraph{Sampling and tree reduction.}
We consider a three-stage horizon with decision times
$T_1=\Delta t, T_2=2\Delta t$, and $T_3=3\Delta t$, where $\Delta t=0.5$ (years).
For each defect, the fitted MJP gives a TPM over one decision interval, cf. \cref{sec:transient_dist}. In this example, $\bm A^{(s)}$ is obtained from an MJP fitted with the generator parameterization $\bm{A}_{\RN{3}}$ and the regression setting in \cref{eq:trans_rates_covariates}.
Using the one-step TPMs, we generate Monte Carlo degradation paths on $(T_1,T_2,T_3)$ for a set of $10$ selected defects (sampling each defect independently conditional on initial distributions at $T_1$).
The full joint scenario tree becomes too large to use directly, so we approximate it by a reduced scenario tree using stagewise clustering (weighted K-medoids, cf.\ \citet{park2009a}): first, sampled joint states at $T_2$ are clustered into $K_2$ representative nodes; then, within each stage-2 cluster, the descendant joint states at $T_3$ are clustered into $K_3$ representative children.
Node probabilities are estimated by relative frequencies of the paths, and each node stores the per-defect marginal distributions. These are used to compute (nodewise) the quantities of $q^{(s)}(t)$ and $\bar c^{(s)}(t)$.

\paragraph{Encoding of endogenous information.}
The reduced tree represents future joint defect scenarios, while inspection changes what is known later.
If defect $s$ is not inspected by stage $T_t$, then quantities such as $q^{(s)}(T_t)$ and $\bar c^{(s)}(T_t)$ are computed from the per-defect marginal state distribution stored at that node.
If defect $s$ is inspected and class $j\in\mathcal J$ is observed, then later-stage $q^{(s)}(t)$ and $\bar c^{(s)}(t)$ are computed conditional on that observation.
In this way, the reduced scenario tree stays fixed, but the quantities entering the optimization model switch to observation-conditional values when inspection occurs.


\paragraph{Numerical details.}
The reduced-tree architecture is specified by $(K_2,K_3)$, where $K_2$ is the number of stage-2 clusters and $K_3$ the number of stage-3 children per stage-2 cluster.
Candidate architectures are evaluated by constructing the reduced tree from a training set of Monte Carlo paths, solving the resulting mixed-integer program, and estimating the induced policy’s mean total cost on an independent validation set. We select the architecture with the lowest validation mean total cost (with smaller-tree tie-breaks).

\Cref{tab:setup} displays the parameters used to run the optimization program.
\begin{table}[htb!]
\centering
\caption{Experimental setup.} \label{tab:setup}
\begin{tabular}{p{0.43\linewidth}p{0.47\linewidth}}
\toprule
Item & Value \\
\midrule
Number of defects & \(10\) \\
Decision times & \((0.5,\,1.0,\,1.5)\) \\
Training / validation / test paths & \(10^6 / 10^5 / 10^5\) \\
Training / validation / test seeds & \(42 / 43 / 44\) \\
Inspection cost \(c_{\mathrm{ins}}\) & \(150\) \\
Failure cost \(c_{\mathrm{fail}}\) & \(200{,}000\) \\
Maintenance costs for classes in $\mathcal E$ & \(2000,\,3500,\,4500,\,8000,\,10000\) \\
Stage budgets & \((20000,\,19500,\,18000)\) \\
Critical classes $\mathcal E_{\text{critical}}$ & \(\{ \text{Class } $1$, \text{Class } $0$\}\) \\
\bottomrule
\end{tabular}
\end{table}

\Cref{tab:reduced_tree} displays the parameters used to reduce the full tree and obtained results from solving the program.
\begin{table}[htb!]
\centering
\caption{Summary for the selected reduced-tree architecture.} \label{tab:reduced_tree}
\begin{tabular}{p{0.55\linewidth}p{0.30\linewidth}}
\toprule
Item & Value \\
\midrule
Selected architecture $(K_2 \times K_3)$ & $20\times 8$ \\
Realized stage-2 / stage-3 nodes & $20 / 160$ \\
Validation mean total cost (s.e.) & $105063.69\ (364.86)$ \\
Test mean total cost & $104854.52$ \\
Stage-3 critical-count total variation & $6.9\times 10^{-5}$ \\
Stage-1 budget usage (\%) & $96.64\%$ \\
\bottomrule
\end{tabular}
\end{table}








